% RustGuard --- draft paper.
%
% Target: IACR Transactions on Cryptographic Hardware and Embedded Systems.
% To switch to the TCHES class, replace the documentclass line with
%   \documentclass[journal=tches,submission,spdx=false]{iacrtrans}
% and drop the geometry/titling lines below; everything else is portable.
%
% BEFORE SUBMISSION: every citation in refs.bib is from memory and MUST be
% checked against the actual paper (authors, venue, year, and above all the
% claim attributed to it). Do not submit on trust.
\documentclass[11pt,a4paper]{article}
\usepackage[margin=1in]{geometry}
\usepackage[T1]{fontenc}
\usepackage{booktabs}
\usepackage{graphicx}
\usepackage{amsmath}
\usepackage{url}
\usepackage[hidelinks]{hyperref}
\usepackage{xcolor}

\newcommand{\todo}[1]{\textcolor{red}{[TODO: #1]}}

\title{Does Constant-Time Rust Stay Constant-Time on Embedded Silicon?\\
       A Systematic, Equipment-Free Evaluation}
\author{Taha Hunaid Ali}
\date{Draft --- \today}

\begin{document}
\maketitle

\begin{abstract}
Memory-safe languages are now recommended for security-critical firmware, and
the Rust cryptographic ecosystem advertises constant-time execution through the
\texttt{subtle} crate and branch-free implementation strategies. Whether those
source-level guarantees survive compilation and reach the silicon that actually
ships is an empirical question, and one that is usually answered with equipment
most developers do not have.

We report a systematic constant-time evaluation of eleven symmetric and two
public-key Rust implementations on two Cortex-M4 microcontrollers from different
vendors, using nothing but the core's own cycle counter. The study spans four
optimization levels, five core-frequency configurations that change flash wait
states, and two experiment designs that separate a primitive's comparison path
from its arithmetic core: 257 captures and roughly 750\,000 on-chip
measurements in total.

Every production implementation is cycle-invariant in every configuration. We
make that null result quantitative by injecting calibrated leaks of known size:
on this platform a mean difference of $0.49$ cycles is detected, so ``no leakage
observed'' becomes ``any leak above half a cycle would have been observed''. We
further report three findings that the grid makes visible. First, LLVM
\emph{removes} a deliberate early-exit timing leak at \texttt{-O2} and above ---
identically on both vendors' silicon --- which means a positive control
validated at one optimization level does not validate the others, and studies
that validate once may be reporting unvalidated columns. Second, flash wait
states inflate cycle counts by between $4.3\%$ and $27.2\%$ depending on whether
a primitive is fetch- or compute-bound, without introducing data-dependent
timing. Third, binary-level branch-counting triage produces eight false
positives out of ten flagged primitives against measured ground truth, which
bounds what such screening can claim.
\end{abstract}

\section{Introduction}

Constant-time implementation is the standard defence against timing attacks on
cryptographic code, and it is a property of the \emph{binary}, not of the
source. Compilers are permitted to rewrite branch-free source into branching
machine code, and a growing literature documents cases where they do. Rust adds
a new dimension to this old problem: its cryptographic ecosystem is young,
widely deployed in embedded contexts, and built on the premise that safety
properties expressed in the type system and in libraries such as \texttt{subtle}
survive to the artefact.

Most constant-time evaluation on embedded targets assumes laboratory equipment
--- an oscilloscope, a ChipWhisperer-class capture rig, or a trace probe. That
assumption shapes which questions get asked and by whom. We take the opposite
constraint as a design goal: every measurement in this paper is taken with the
Cortex-M's own DWT cycle counter on a development board costing under
\$20, and the entire study is reproducible by anyone holding the same board.

\paragraph{Contributions.}
\begin{itemize}
  \item A systematic constant-time evaluation of thirteen Rust cryptographic
        implementations --- nine AEADs and MACs, two public-key primitives, and
        two controls --- across two silicon vendors, four optimization levels
        and five core-frequency configurations (Section~\ref{sec:results}).
  \item A \emph{calibrated leak ladder} that converts a null result into a
        bounded one by measuring the smallest leak the platform can detect
        (Section~\ref{sec:floor}).
  \item The observation that a realistic positive control is itself
        optimizer-dependent: the early-exit comparison we use as a control stops
        leaking at \texttt{-O2}, on both vendors, so per-configuration control
        validation is necessary rather than good practice
        (Section~\ref{sec:optimizer}).
  \item A quantified comparison of binary-level screening against measured
        ground truth on the same image (Section~\ref{sec:static}).
  \item An open artefact: the probe registry, both firmware harnesses, the
        capture drivers and the analysis pipeline, with every raw capture.
\end{itemize}

\section{Background}

\paragraph{Constant-time code and its enemies.}
An implementation is constant-time when its execution time is independent of
secret data. \todo{cite Kocher; Bernstein cache-timing; Almeida et al.\ ct-verif}

\paragraph{dudect.}
Rather than model the microarchitecture, dudect~\cite{dudect} tests the
hypothesis that two input classes are timed identically, using Welch's
$t$-test with a $|t| > 4.5$ threshold. We adopt its design, including
percentile cropping for noisy platforms.

\paragraph{What can leak on a Cortex-M4.}
The M4 has no data cache and no branch predictor, so the leak classes that
dominate the literature on application processors --- cache-line-granular table
lookups, speculative execution --- cannot manifest. What remains is
(i)~secret-dependent branching, (ii)~variable-latency instructions, notably
\texttt{UDIV}/\texttt{SDIV} at 2--12 cycles, and (iii)~early-return comparison.
This narrows the threat model in a way that makes an exhaustive, equipment-free
study feasible, and it is why our conclusions are stated for this class of core
rather than for Rust cryptography in general.

\section{Methodology}

\subsection{A uniform probe registry}
Each implementation under test is wrapped as a \emph{probe} exposing the same
three operations: produce a genuine tag, verify a supplied tag, and authenticate
a fixed message under a supplied key. The firmware carries every probe in a
single image and selects between them over a serial protocol, so a full sweep is
one flash per configuration rather than one per crate, and adding an
implementation to the study is one registry entry.

\subsection{Two experiments}
\label{sec:experiments}
A single experiment design would have overstated our results, so we run two.

\paragraph{\textsc{verify} --- the comparison path.}
The key is fixed and the tag varies. The fixed class presents the genuine tag
with its last byte flipped, so a byte-wise comparison runs nearly to completion;
the random class presents a uniformly random tag, which mismatches almost
immediately. Both classes are rejected, so both execute the same failure path
and only the comparison differs.

\paragraph{\textsc{keyed} --- the arithmetic core.}
The classic dudect design: a fixed key against uniformly random keys, timing the
implementation authenticating a fixed message. This is the experiment that
exercises the AES key schedule and S-box, the GHASH and Poly1305 field
arithmetic, the permutation, and --- for the public-key probes --- scalar
multiplication. A clean \textsc{verify} verdict says nothing about the core, and
we report the two matrices separately for that reason.

\subsection{Controls, per configuration}
\label{sec:controls}
A null result is worthless without a positive control, and --- as
Section~\ref{sec:optimizer} shows --- a control validated in one configuration
may be silently inert in another. We therefore require every column of every
matrix to contain a control that demonstrably leaks, and we use two:

\begin{description}
  \item[\textsc{leaky}] a realistic leak: an early-exit tag comparison, the
        canonical mistake this kind of study exists to catch.
  \item[\textsc{canary}] an optimizer-proof leak: work proportional to a secret
        byte, behind \texttt{core::hint::black\_box}, which the compiler may not
        reason through. It leaks by construction at every optimization level.
\end{description}

\subsection{Measurement}
Measurements are taken with the DWT cycle counter with interrupts masked, and
the secret input is passed through \texttt{black\_box} so the operation cannot
be hoisted out of the measured region. On an in-order core with no cache and no
preemption, the resulting count is a deterministic function of the input: the
same input always yields exactly the same number. This is a much stronger
measurement regime than an OS-scheduled core affords, and it is what allows the
zero-spread and detection-floor claims below.

\section{Experimental setup}

\begin{table}[t]
\centering
\caption{Measurement grid. All configurations carry a validated control.}
\label{tab:setup}
\begin{tabular}{ll}
\toprule
Axis & Values \\
\midrule
Silicon      & TM4C123GH6PM (TI), STM32F303VC (ST) --- both Cortex-M4 \\
Core clock   & 16, 80\,MHz (TM4C); 8, 32, 64\,MHz (STM32) \\
Flash latency& 0--2 wait states, following the clock \\
Optimization & \texttt{-O0}, \texttt{-O1}, \texttt{-O2}, \texttt{-O3} \\
Experiments  & \textsc{verify}, \textsc{keyed} \\
Primitives   & 9 AEAD/MAC, 2 public-key, 2 controls, 7 ladder rungs \\
Traces       & 3000 per symmetric cell; 400 per public-key cell \\
Total        & 257 captures, ${\approx}750\,200$ measurements \\
\bottomrule
\end{tabular}
\end{table}

Table~\ref{tab:setup} summarises the grid. The implementations evaluated are
\texttt{aes-gcm}, \texttt{aes-gcm-siv}, \texttt{chacha20poly1305},
\texttt{ascon-aead}, \texttt{eax}, \texttt{ccm}, \texttt{hmac-sha256},
\texttt{cmac-aes128}, an in-house ASCON-128, and the public-key primitives
\texttt{x25519-dalek} and \texttt{p256} scalar multiplication. Version numbers
are pinned in the artefact.

\section{Results}
\label{sec:results}

\subsection{Symmetric primitives are cycle-invariant}

\begin{table}[t]
\centering
\caption{Verification cost on the TM4C123 at \texttt{-O3}, 16\,MHz. Every
implementation showed a cycle-count spread of exactly zero across all traces in
both experiments and in every configuration.}
\label{tab:cycles}
\begin{tabular}{lrr}
\toprule
Implementation & Cycles (\textsc{verify}) & Spread \\
\midrule
\texttt{chacha20poly1305} & 7\,982 & 0 \\
\texttt{rustguard-ascon128} & 8\,217 & 0 \\
\texttt{ascon-aead} & 8\,769 & 0 \\
\texttt{hmac-sha256} & 12\,772 & 0 \\
\texttt{cmac-aes128} & 16\,445 & 0 \\
\texttt{aes-gcm} & 29\,589 & 0 \\
\texttt{ccm-aes128} & 41\,342 & 0 \\
\texttt{aes-gcm-siv} & 78\,794 & 0 \\
\texttt{eax-aes128} & 121\,707 & 0 \\
\bottomrule
\end{tabular}
\end{table}

Across the whole grid, no production implementation was flagged in any
configuration. The stronger statement is the one in Table~\ref{tab:cycles}:
in 108 of 130 \textsc{verify} captures and 116 of 127 \textsc{keyed} captures,
the observed cycle count was \emph{identical for every trace} --- a spread of
zero. The remaining captures are exactly the controls and ladder rungs. These
implementations do not merely fall below a detection threshold; over the sampled
input space they exhibit no input-dependent timing at all.

\subsection{Public-key primitives}

Symmetric primitives are the easy case: branch-free by construction. Scalar
multiplication is where data-dependent branching and modular arithmetic
traditionally cause trouble, and to our knowledge these crates have not
previously been measured on embedded silicon. \texttt{x25519-dalek} costs
$1\,963\,173$ cycles and \texttt{p256} scalar multiplication $6\,689\,177$
cycles on the TM4C at \texttt{-O3}; both showed a spread of zero across random
scalars, in both experiments and at every optimization level measured.

\subsection{The optimizer removes a real timing leak}
\label{sec:optimizer}

\begin{table}[t]
\centering
\caption{The realistic positive control across optimization levels. LLVM
rewrites the early-exit comparison into branchless code at \texttt{-O2}, and the
control stops leaking --- on both vendors. The optimizer-proof canary leaks
everywhere.}
\label{tab:optimizer}
\begin{tabular}{lrrrr}
\toprule
 & \multicolumn{2}{c}{TM4C123} & \multicolumn{2}{c}{STM32F303} \\
\cmidrule(lr){2-3}\cmidrule(lr){4-5}
Level & \textsc{leaky} $|t|$ & \textsc{canary} $|t|$ & \textsc{leaky} $|t|$ & \textsc{canary} $|t|$ \\
\midrule
\texttt{-O0} & 12\,997 & 61.28 & $\infty$ & 61.55 \\
\texttt{-O1} & 11\,259 & 61.28 & 11\,259 & 61.55 \\
\texttt{-O2} & 0.00 & 61.28 & 0.00 & 61.55 \\
\texttt{-O3} & 0.00 & 61.28 & 0.00 & 61.55 \\
\bottomrule
\end{tabular}
\end{table}

Table~\ref{tab:optimizer} reports the finding we consider most useful to other
practitioners. The early-exit comparison leaks enormously at \texttt{-O0} and
\texttt{-O1} and then, at \texttt{-O2}, stops leaking entirely: LLVM rewrites it
into branchless code. The effect reproduces identically on both vendors'
silicon, which locates it in the compiler rather than in a microarchitecture.

The methodological consequence matters more than the curiosity. A study that
validates its harness once --- at one optimization level, as is common --- and
then reports clean results at other levels may be reporting columns in which its
control was inert, and therefore claiming detection capability it did not have
at the time. This is not hypothetical: our own \texttt{-O2} and \texttt{-O3}
columns were briefly in exactly that state before the canary was added.

A second consequence is more cheerful: compilers are usually cast as the
adversary of constant-time code, and here LLVM repaired a genuine leak. We do
not suggest relying on this. It is an optimization decision, not a guarantee,
and the same mechanism that removed the leak at \texttt{-O2} could reintroduce
one elsewhere.

\subsection{Flash wait states}
\label{sec:clock}

Raising the core clock past the flash's sustainable rate forces wait states.
The same binary then costs more cycles --- but how many more depends on the
primitive. ASCON verification inflates by $27.2\%$ from 16 to 80\,MHz on the
TM4C and by $25.9\%$ from 8 to 32\,MHz on the STM32, while
\texttt{hmac-sha256} inflates by only $4.3\%$ and \texttt{aes-gcm} by $6.4\%$:
fetch-bound code pays for wait states, compute-bound code hides them behind
work it was doing anyway. No configuration introduced data-dependent timing.

\subsection{Binary screening against measured ground truth}
\label{sec:static}

Counting conditional branches in a disassembly is a cheap way to triage a
dependency tree, and it is tempting to treat the result as evidence. Joining our
census against measurements from the same image gives the error rate directly:
of ten primitives flagged as carrying potentially data-dependent constructs,
\emph{eight are false positives}, and none of the genuine leaks was missed.
\texttt{chacha20poly1305} carries 45 conditional branches and is perfectly
cycle-invariant --- the branches are driven by block counts, not secrets.
Branch-counting prioritises what to measure; it does not detect leaks.

\subsection{How small a leak can we see?}
\label{sec:floor}

\begin{table}[t]
\centering
\caption{Calibrated leak ladder on the TM4C123 at \texttt{-O3}. The smallest
rung still flagged sets the platform's detection floor.}
\label{tab:ladder}
\begin{tabular}{lrrr}
\toprule
Injected leak & Mean difference (cycles) & $|t|$ & Detected \\
\midrule
1 iteration    & 0.49 & 37.60 & yes \\
2 iterations   & 2.91 & 37.60 & yes \\
4 iterations   & 7.28 & 37.60 & yes \\
8 iterations   & 17.47 & 37.60 & yes \\
16 iterations  & 62.12 & 37.60 & yes \\
64 iterations  & 248.01 & 37.60 & yes \\
256 iterations & 994.45 & 37.60 & yes \\
\bottomrule
\end{tabular}
\end{table}

Every clean verdict above is bounded by Table~\ref{tab:ladder}: the smallest
leak we injected, a single extra loop iteration producing a mean difference of
$0.49$ cycles, is detected. Sub-cycle sensitivity is possible because the
counter is exact and the injected leak makes the distribution bimodal, so the
\emph{mean} difference can be a fraction of a cycle while every individual
measurement is an integer.

Two properties of the ladder deserve comment. First, $|t|$ is identical at every
rung, because Welch's $t$ is invariant under affine scaling of the measurements:
enlarging the leak scales the mean difference and the standard deviation
together. The statistic therefore reports \emph{whether} the classes separate,
not by how much, and a larger leak is not a stronger signal in this regime.
Second, and consequently, the ladder establishes an upper bound on the floor
rather than the floor itself: detection does not degrade as the leak shrinks to
one iteration, so the true floor lies at or below $0.49$ cycles of mean
difference, and would need a sub-iteration leak to locate exactly. We therefore
state our results as: any leak above half a cycle of mean difference would have
been detected, and none was. On a platform where measurement is noisy this same
ladder yields a floor orders of magnitude higher, which is the comparison we
intend it for.

\section{Discussion and limitations}

\paragraph{What this study does not cover.}
We measure \emph{timing}, not power or electromagnetic emission; constant-time
code can still leak through those channels, and a capture rig is required to say
otherwise. Our conclusions are stated for in-order, cacheless Cortex-M4 cores:
the leak classes we exclude by architecture --- cache-line timing, speculation
--- are precisely those that dominate on application processors, and a
constant-time verdict here does not transfer to those targets. Extending the
same harness to application-class ARM cores is implemented and pending hardware
\todo{fold in Raspberry Pi 3/4/5 results}.

\paragraph{Sampled, not exhaustive.}
Zero spread is observed over the sampled input space, not proven over all
inputs. It is a strong empirical statement, not a proof, and it complements
rather than replaces static verification.

\paragraph{Two probe defects we found by checking.}
Both are worth recording because they are the kind of error that produces a
confident, wrong result. Our \texttt{x25519} probe initially compared a derived
public key with a plain slice comparison, which is an early-exit memcmp; it
measured a 279-cycle difference that we would have been entitled to report as a
leak in a library whose entire purpose is constant-time execution. The leak was
in our harness. Separately, the \texttt{p256} probe declared a 33-byte
compressed point against a 32-byte buffer and silently failed its verify
captures. Neither would have been caught by a pipeline that trusted its own
outputs.

\section{Related work}
\todo{Write properly, and verify every citation against the source: dudect
(Reparaz et al., DATE 2017); compiler-induced violations (Kaufmann et al., CANS
2016; Simon, Chisnall \& Anderson, EuroS\&P 2018); static/symbolic CT analysis
(ct-verif, Binsec/Rel); developer practice (Jancar et al., IEEE S\&P 2022);
MicroWalk. Position this work as: systematic, multi-implementation,
equipment-free, and quantified by a measured detection floor.}

\section{Conclusion}

Across 257 captures on two vendors' Cortex-M4 silicon, thirteen Rust
cryptographic implementations --- including two public-key primitives --- showed
no input-dependent timing, at a measured sensitivity of half a cycle. The
ecosystem's constant-time claims hold on this class of target. The results that
travel beyond this ecosystem are methodological: positive controls are
optimization-dependent and must be validated per configuration, binary
branch-counting over-approximates by a factor we can now quote, and a null
timing result should be published with the detection floor that bounds it.

\bibliographystyle{plain}
\bibliography{refs}

\end{document}
